\documentclass[10pt,a4paper]{amsart}
\usepackage[margin=19mm]{geometry}
\usepackage{amsmath,amssymb,mathtools}
\usepackage[scr=rsfs,cal=euler]{mathalfa}
\usepackage{xfrac}
\usepackage[unicode,hidelinks,bookmarks=false]{hyperref}
\newcommand{\ie}{i.e.\ }
\newcommand{\cf}{cf.\ }
\newcommand{\resp}{respectively\ }
\newcommand{\Subsection}{\subsection}
\providecommand{\nobreakdash}{\nobreak\hbox{-}}
\setlength{\emergencystretch}{2em}
\multlinegap0pt
\usepackage{bm}
\usepackage{bbm}
\usepackage{dsfont}
\usepackage{xspace}
\usepackage{cancel}
\usepackage{mathtools}
\usepackage{amsthm}
\makeatletter
\@ifundefined{theorem}{\newtheorem{theorem}{Theorem}}{}
\@ifundefined{proposition}{\newtheorem{proposition}[theorem]{Proposition}}{}
\makeatother
\newcommand{\Eul}{\mathcal{E}}
\newcommand{\Pj}{\mathbf{P}}
\newcommand{\Q}{\mathbf{Q}}
\title[Quartic Rational Diophantine Quadruples and the Euler Surfac]{Quartic Rational Diophantine Quadruples and the Euler Surface}
\author{Alen Andra\v{s}ek, Matija Kazalicki, Domagoj Vlah}
\date{Source: arXiv:2604.19140v1 (CC BY 4.0)}
\begin{document}
\maketitle

\noindent\emph{Source and licence.} Quartic Rational Diophantine Quadruples and the Euler Surface. Version arXiv:2604.19140v1 (CC BY 4.0), distributed under the Creative Commons Attribution 4.0 International licence, as recorded in the arXiv licence metadata for this exact version and shown on the abstract page https://arxiv.org/abs/2604.19140v1. Full rights notice: licenses/arxiv-2604.19140-CC-BY-4.0.txt.\par\smallskip

\noindent\textit{Editorial scope.} Counted unit: the Proposition whose source label is \texttt{prop:euler-to-quadruple} in the article (source0.tex lines 105--133, source label \texttt{prop:euler-to-quadruple}), delivered together with the article's own complete original proof of it (source0.tex lines 417--457, one \texttt{proof} environment of 41 lines). No mathematics is written, repaired or re-derived: statement and proof are the article's own text line for line. The proof is detached from the statement in the source: the article states the result at lines 105--133 and prints its proof at lines 417--457, and the proof environment's own header names the statement, so the association is the article's own. Required context retained uncounted: eq:abcd-from-rstuvw (lines 195--200); eq:r-t-from-k (lines 301--303); eq:def-s-special (lines 305--307); eq:special-quartic (lines 312--314); prop:auto-square-special (lines 348--353); prop:euler-specialization (lines 377--395). Every definition, display and label of those lines is retained unchanged. Omitted from this package and not redistributed: 1--104, 134--194, 201--300, 304--304, 308--311, 315--347, 354--376, 396--416, 458--767. Nothing inside a retained excerpt is omitted or altered: every retained body is delivered exactly as the article's lines are written, so no omission receipt of any kind accompanies this package. Citations: the retained excerpts carry no citation command at all, so no bibliography entry of the article is transcribed and no reference is redistributed. Reference adaptation: none was needed and none was made; every reference inside the delivered unit resolves inside it, so no unresolved reference or citation is produced. Printed slips kept literal: no printed word, symbol, hypothesis or number of the counted statement or of its proof was altered. Layout and class: the article's own class is \texttt{article} with packages including geometry, fontenc, inputenc, lmodern, microtype, amsmath, amssymb, amsthm, mathtools, hyperref, enumitem, mathrsfs; the wrapper is the lane's pinned \texttt{amsart} editorial wrapper at 10pt on A4, which restates the theorem-like environments the retained text needs and adds the article's own macro definitions that the retained text needs (\texttt{\textbackslash Eul}, \texttt{\textbackslash Pj}, \texttt{\textbackslash Q}), with extra wrapper packages (bm, bbm, dsfont, xspace, cancel, mathtools). The delivered numbering is the unit's own. Assets: no retained span contains a figure environment, a graphics-inclusion command, a PDF-page inclusion, a table or a TikZ picture; no image file is copied into this package, the package declares no asset dependency and no image file is redistributed with it. Licence: version arXiv:2604.19140v1 of this article is distributed under the Creative Commons Attribution 4.0 International licence, as recorded in the arXiv licence metadata for this exact version and shown on the abstract page https://arxiv.org/abs/2604.19140v1. A full rights notice is delivered at licenses/arxiv-2604.19140-CC-BY-4.0.txt. Characters: every retained excerpt is pure ASCII, so the delivered engine, XeLaTeX with its default Latin Modern text font, reports no missing character.
\begin{equation}\label{eq:abcd-from-rstuvw}
a^2=\frac{(r^4-1)(u^4-1)}{v^4-1},\qquad
b=\frac{r^4-1}{a},\qquad
c=\frac{s^4-1}{a},\qquad
d=\frac{u^4-1}{a},
\end{equation}

\begin{equation}\label{eq:r-t-from-k}
r=\kappa w,\qquad t=\kappa u.
\end{equation}

\begin{equation}\label{eq:def-s-special}
s=\kappa uw.
\end{equation}

\begin{equation}\label{eq:special-quartic}
u^4+w^4=1+\frac{1}{\kappa^4}.
\end{equation}

\begin{proposition}\label{prop:auto-square-special}
Let $\kappa,u,w\in\Q^\times$ and define $r,t,s,v$ by \eqref{eq:r-t-from-k}, \eqref{eq:def-s-special}, and $v=0$. If \eqref{eq:special-quartic} holds, then the square condition in \eqref{eq:abcd-from-rstuvw} is automatic. More precisely,
\[
(r^4-1)(u^4-1)(v^4-1)=\left(\frac{r^4-1}{\kappa^2}\right)^2.
\]
\end{proposition}

\begin{proposition}\label{prop:euler-specialization}
Let $(X:Y:Z:W)\in\Pj^3(\Q)$ with $ZW\neq 0$, and define
\begin{equation}\label{eq:k-u-w-from-euler}
\kappa=\frac{W}{Z},\qquad u=\frac{X}{W},\qquad w=\frac{Y}{W}.
\end{equation}
Then $(X:Y:Z:W)$ lies on the Euler surface
\[
\Eul:\qquad X^4+Y^4=Z^4+W^4
\]
if and only if $k,u,w$ satisfy \eqref{eq:special-quartic}. Equivalently, if we set
\[
r=\frac{Y}{Z},\qquad s=\frac{XY}{ZW},\qquad t=\frac{X}{Z},\qquad v=0,
\]
then
\[
r=\kappa w,\qquad s=\kappa uw,\qquad t=\kappa u,
\]
and the compatibility conditions hold.
\end{proposition}

\begin{proposition}\label{prop:euler-to-quadruple}
Let $(X:Y:Z:W)\in \Eul(\Q)$ with $ZW\neq 0$, and define
\begin{equation}\label{eq:euler-map}
 a=\frac{X^4-W^4}{Z^2W^2},\qquad
 b=-\frac{W^2}{Z^2},\qquad
 c=\frac{Y^4-W^4}{Z^2W^2},\qquad
 d=\frac{Z^2}{W^2}.
\end{equation}
Then
\[
ab+1=\left(\frac{Y}{Z}\right)^4,
\qquad
ac+1=\left(\frac{XY}{ZW}\right)^4,
\qquad
ad+1=\left(\frac{X}{W}\right)^4,
\]
\[
bc+1=\left(\frac{X}{Z}\right)^4,
\qquad
bd+1=0,
\qquad
cd+1=\left(\frac{Y}{W}\right)^4.
\]
Moreover, the four numbers in \eqref{eq:euler-map} are pairwise distinct and nonzero if and only if
\begin{equation}\label{eq:nondeg-euler}
XYZW\,(X^2-Y^2)(X^2-W^2)(Y^2-W^2)\neq 0.
\end{equation}
In that case they form a quartic rational Diophantine quadruple.
\end{proposition}

\begin{proof}[Proof of Proposition~\ref{prop:euler-to-quadruple}]
Let $(X:Y:Z:W)\in\Eul(\Q)$ with $ZW\neq 0$, and define $\kappa,u,w,r,s,t,v$ as in Proposition~\ref{prop:euler-specialization}. Then $v=0$, $r=\kappa w$, $s=\kappa uw$, $t=\kappa u$, and \eqref{eq:special-quartic} holds. By Proposition~\ref{prop:auto-square-special}, the square condition in \eqref{eq:abcd-from-rstuvw} is automatic, and therefore the formulas \eqref{eq:abcd-from-rstuvw} produce a quartic quadruple.

Substituting
\[
r=\frac{Y}{Z},\qquad s=\frac{XY}{ZW},\qquad u=\frac{X}{W},\qquad v=0
\]
into \eqref{eq:abcd-from-rstuvw}, and choosing the sign of $a$ so that
\[
a=\frac{X^4-W^4}{Z^2W^2},
\]
we obtain exactly the formulas \eqref{eq:euler-map}. The displayed identities for $ab+1$, $ac+1$, $ad+1$, $bc+1$, $bd+1$, and $cd+1$ are then immediate.

It remains to determine when $a,b,c,d$ are pairwise distinct and nonzero. Since $ZW\neq 0$, the numbers $b$ and $d$ are automatically nonzero. Also,
\[
a=0 \iff X^4=W^4 \iff X^2=W^2,
\qquad
c=0 \iff Y^4=W^4 \iff Y^2=W^2.
\]
Further,
\[
a-b=\frac{X^4}{Z^2W^2},
\qquad
c-d=-\frac{X^4}{Z^2W^2},
\]
so $a=b$ and $c=d$ are both equivalent to $X=0$. Likewise,
\[
a-d=-\frac{Y^4}{Z^2W^2},
\qquad
c-b=\frac{Y^4}{Z^2W^2},
\]
using $X^4+Y^4=Z^4+W^4$, so $a=d$ and $c=b$ are both equivalent to $Y=0$. Finally,
\[
a-c=\frac{X^4-Y^4}{Z^2W^2},
\]
so $a=c$ is equivalent to $X^2=Y^2$. Since $bd+1=0$, we also have $bd=-1$, and therefore $b\neq d$. This proves that $a,b,c,d$ are pairwise distinct and nonzero if and only if
\[
XYZW(X^2-Y^2)(X^2-W^2)(Y^2-W^2)\neq 0,
\]
which is exactly the condition stated in Proposition~\ref{prop:euler-to-quadruple}.
\end{proof}

\end{document}
